\chapter{Future Work}\label{chap:future_work}

The next steps should focus on turning the verified workflow into a validated
numerical forming limit campaign.

\begin{enumerate}
  \item \textbf{Complete the mesh verification.} The current
  $\Delta t_\mathrm{MS}=\SI{1e-6}{\second}$ and $f_\mathrm{MR}=1.41$ setting is
  only a working choice for the \texttt{W100} reference case. The missing
  $f_\mathrm{MR}=1.5$ case should be completed, and the mesh study should be
  repeated for the relevant thicknesses and punch configurations before mesh
  convergence is claimed.

  \item \textbf{Run the numerical campaign.} The simulations should be extended
  from the reference cases to the specimen widths, sheet thicknesses, punch
  diameters, and material orientations used experimentally. This is required
  before a numerical FLC can be compared with the experimental campaign.

  \item \textbf{Validate material and contact parameters.} The AA5182
  surrogate material constants should be replaced by an AA5754 calibration once
  the required characterisation tests are available. The punch and die friction
  coefficients should also be calibrated against experimental
  force displacement curves, since the local friction check showed a strong
  influence on the fracture location.

  \item \textbf{Compare forming limit identification criteria.} The simulated
  strain histories should be processed with the same ISO, Volk Hora, and
  Min Stoughton logic used for the experiments. This would separate numerical
  model effects from post processing effects.

  \item \textbf{Improve usability.} The Streamlit interface should be kept as a
  practical front end for job setup, result synchronisation, and plot review,
  with emphasis on faster result loading and clearer status feedback.
\end{enumerate}
